\include{zLaTex-Preamble}

\title{Project 1 - Design and Test Writeup} 
\author{Joshua Miller}
%% \date{} **/

\begin{document}
\rhead{\fancyplain{}{Josh Miller}}
\maketitle



\section{General Goals}

The goal for this project is to optimize the performance of
Floyd-Warshall algorithm by reducing execution time using
pthreads. 

\section{Serial Design}

\subsection{Serial Optimization}

A naive implementation of Floyd's algorithm consists of a simple
triple loop. In order to begin the parallelization with efficient
code, I plan to test a version of the algorithm using an additional
transpose adjacency matrix.  Given matrix $A$, transposed matrix $A^T$
will allow for more optimal use of the cache by exploiting the memory
locality.  Floyd's algorithm can best be implemented with
$\mathit{O}(n^3)$ while the transposition of the matrix executes with
$\mathit{O}(n^2)$ complexity. Thus, the overhead presented by this
will be minimal in comparison with the operations required for the
algorithm itself for large problem sizes.  However, I expect the
negative impact in performance to be observed in the execution time of
the transposed implementation on smaller problem sizes. The percentage
of overhead to be introduced by the transposition would theoretically
follow as 
\al{
        \text{overhead } = \frac1{n^3/n^2} = \frac{1}{n}
}

Therefore for problem sizes of $n \geq 100$, the overhead introduced
will be 1\%.

Because both the original and transposed version of the adjacency
matrix must be updated with each iteration, it will require an extra
memory write per matrix location.  However, I predict that the benefit
from the introduced locality will still reduce the execution time,
even with the included overhead from the transposition.

An additional optimization from the naive implementation is presented
in the boundary conditions of the adjacency matrix.  We know that any
edge weight of 10000000 \emph{(INFINITY)} represents two vertices
without a direct path and are thus able to skip the addition and
comparison that would have been required for each matrix location
using that value.

\lang{C}
\begin{code}

/* reads a matrix in from a file returns with width N.
 * sets int *m to be location of newly read in matrix */
int fread_matrix(int fd, int *m);

/* allocates space for new matrix array and returns transposition of m1 */
int * transpose_matrix(int * m1, int n); 

/* takes matrix at m with width n and performs tripple nested loop */
int * floyds(int *m, int n);

\end{code}

Because I do not know whether the matrix transformed version will
actually be more efficient I will compare the two implementations.  I
will describe the threaded design as if I were using the transposed
implementation, but will implement the same design without the
transposed matrix for comparison as well.


\section{Threaded Design}


Considering the original adjacency matrix $A$ and a transposed $A^T$, 

\al{
A = 
  \m{
     a_{00} & a_{10} & a_{20} & a_{30} \\
     a_{01} & a_{11} & a_{21} & a_{31} \\
     a_{02} & a_{12} & a_{22} & a_{32} \\
     a_{03} & a_{13} & a_{23} & a_{33} 
  },~
A^T = 
  \m{
     a_{00} & a_{01} & a_{02} & a_{03} \\
     a_{10} & a_{11} & a_{12} & a_{13} \\
     a_{20} & a_{21} & a_{22} & a_{23} \\
     a_{30} & a_{31} & a_{32} & a_{33} 
  }
}

We can divide the work of the algorithm at the level of the second
nested loop.  I am unsure whether there exists a solution to the
division of work by dividing the outermost loop amongst threads,
however I hypothesize that the division in the second loop will
balance the complexity and load balance efficiently.  Be dividing the
work on the innermost level of the nest, we are passing but a few
instructions to each thread. I believe this division will create more
overhead than is productive.  

Therefore, the threaded version of the algorithm will work as follows.  
\begin{enumerate}
\item read in adjacency matrix

\item initialize N\_THREADS threads (cli specified)

\item create all threads 

\item let each thread check out a range of threads (cli specified)

\item wait until all threads report complete

\item increment \tt{k} and repeat from step 4 while \tt{k < n}, broadcast condition variable

\end{enumerate}

Each thread will follow these steps
\begin{enumerate}


\item check out range of rows on mutex protected \emph{current\_row} using \emph{checkout\_rows}

\item skipping row $k$, and ignoring any cases where $A[i,j] = 0$, or $A[i,:] =$ \emph{INFINITY} perform the second level loop of the algorithm, stepping along the row in the threads checked out range.

\item report that thread has completed

\item attempt to checkout new range of rows from \emph{checkout\_rows}. If the function returns a negative value, wait on pthread\_cond\_wait for main thread to signal new $k$ iteration (this will be in a while loop to verify a correct wake-up)


\end{enumerate}

The threads will report completion by turning on a bit coorrelating to
its tag.  The main loop over $k$ will know to increment $k$ when all
threads have reported their completion. Because the test machine only
has 16 cores, a 32 bit integer will suffice.


The definitions from the sequential implementation will remain
unchanged.  The definitions of the functions to be used in the
algorithm are as follows.

\begin{code}

/** condition variable to block worker threads **/
pthread_cond_t k_loop_ready;

/** denotes the next row available to threads **/
volatile int current_row;

/** denotes the next row available to threads **/
volatile int completion_status;

/* number of rows to checkout at a single time for a thread */
int row_bunch;

typedef struct thread_args{
       int n;
       int *matrix;
       int thread_id;
       int row_start;
       int row_end;
} thread_args;

/* handle mutex for current_row{ if current_row > k, return -1
to tell thread to retry later, otherwise update the threads argument's
args->row_start and args->row_end and return 0. } */
int checkout_rows(thread_args *args);

/* handle mutex for completion_status integer{ mask bit offset at
threadid position} */ 
int thread_complete(int threadid);

/* see enumerated flow above */
void * floyd_row(void * thread_args);

\end{code}

I believe that this algorithm will approach an N times speed up in the
region of interest where N is the number of threads assigned. I
believe the overhead will come into play most at the boundary between
the loops over $k$ as well as waiting on the mutex when checking out
rows.  However, with a large matrix, this overhead will be minimal
given that this algorithm lends itself fairly well to parallelization.

\subsection{Resource management}

The following data structures will have wrapper functions that manage
their update in a mutex controlled environment:
\begin{itemize}
\item the current\_row indicator
\item the thread completion indicator
\end{itemize}

Because each thread will be assigned a separate portion of the
matrixm, and should never update outside of that section, these two
mutual exclusions should be sufficient.

\subsection{Communication between threads}

Communication between computation threads will be unnecessary; each
thread will communicate solely to the main process thread by
calling \emph{checkout\_rows} and \emph{thread\_complete}.



\section{Testing}

I plan to create unit tests for both the serial and the parallel
version, though the majority of the algorithm is contained in one
funciton in the serial version. The serial version will be tested by
comparing the output of the algorithm against manually computed
matrices for each step of the process to be read in from a file.
There will be a test for the following conditions as well
\begin{itemize}
\item case of 1$\times$1 matrix
\item case with negative weights, should throw error
\item case with n$\times$m, $n \neq m$ dimensions, should throw error
\item case with infinte lengths (no paths)
\end{itemize}

The parallel version will lend itself better to fine grain testing
because it is split into more detailed functions.  For each value of
$k$, there will be a comparison of the matrix computed with an
expected matrix from a file.  Similarly, each row will be checked for
one matrix. The overall outcome of the parallel version will be
checked against the serial version for random cases.

\subsection{Timing}

THe stopwatch library provided to the class will be used to time the
performance of the implementation to determine how long the thread
took to execute it's computation, how long the thread waited to
checkout its rows, and overall, how long the region of interest took
to execute. 

\subsection{Experiments}

I will perform the following experiments.

\begin{itemize}
\item test the serial performance of transposed vs non-transposed implementations.  I predict that the transposed version will out-perform the non-transposed version for large problem sizes. For explaination see Section~2.1.

\item test the performance of the parallel version with 1 thread.  I predict that the single thread allocated parallel version will execute in a time that is slightly larger than the serial version. But because there will be no overhead introduced by waiting on mutual exclusion flags, the over head will come from the introduced instructions in the parallel version.  Thus I predict that the parallel version will be no larger 5\% different from the serial version.

\item test the performance of the transposed threaded version versus the non-transposed.  I predict that there will be a decrease in execution time for the transposed version.  However, given the symmetry of the algorithm over the i and j indices, this might not even be necessary and may just introduce additional overhead. 

\item test the performance of $2^n$ threads for $n \in \{1,2,3,4\}$ and analyze the way it scales.  I predict that with large matrix sizes, the scaling will approach the idead $2^n$ times. For smaller matrix sizes i.e. $n = 100$, I believe that the time spent handling the mutual exclusion will bve more comparable to the time spent doing the actual computation.

\end{itemize}

\section{Collaboration}

I collaborated with Rachel Hwang on the work load distribution design.

\end{document}
